\chapter{Technische Informationen}
\label{app:TechnischeInfos}

\section{Aktuelle Paketversion}

\begin{center}
	\begin{tabular}{@{}ll@{}}
		\toprule
		Datum    & Datei            \\
		\midrule
		\hgbDate & \texttt{hgb.sty} \\
		\bottomrule
	\end{tabular}
\end{center}


\section{Weitere Details}


Dieses Paket sieht \mbox{UTF-8} kodierte Textdateien vor und
unterstützt \latex\ ausschließlich im direkten PDF-Modus.%
\footnote{Der "klassische" DVI-PS-PDF-Modus von \latex\ wird nicht mehr
unterstützt.}

\subsection{Technische Voraussetzungen}

Eine aktuelle \latex-Installation (\latex-Kern ab Version 2022-06-01,
empfohlen wird jeweils die aktuelle Version von TeX~Live oder MiKTeX) mit
%
\begin{itemize}
		\item \texttt{biber}-Programm (BibTeX-Ersatz) und \texttt{biblatex}-Paket
			in zueinander passenden Versionen,
		\item Latin Modern Schriften (Paket \texttt{lmodern}).%
			\footnote{\url{https://ctan.org/pkg/lm},
				\url{https://tug.org/FontCatalogue/latinmodernroman/}}
\end{itemize}
%
Darüber hinaus ein Texteditor für UTF-8 kodierte (Unicode) Dateien,
sowie Software zum Öffnen und Betrachten von PDF-Dateien.


\subsection{Verwendung unter Windows}
\label{sec:VerwendungUnterWindows}

Eine typische Installation unter Windows sieht folgendermaßen aus:
%
\begin{enumerate}
\item \textbf{MiKTeX}%
	\footnote{\url{https://miktex.org/} -- Achtung:
	Generell wird die \emph{Komplettinstallation} von MiKTeX ("complete
	MiKTeX") empfohlen, da diese bereits alle notwendigen Zusatzpakete und
	Schriftdateien enthält. Bei der Installation ist darauf zu achten,
	dass die automatische Installation erforderlicher Packages durch
	"\emph{Install missing packages on-the-fly: = Yes}" ermöglicht wird
	(NICHT "\emph{Ask me first}")! Außerdem ist zu empfehlen, unmittelbar
	nach der Installation von MiKTeX sowie in weiterer Folge regelmäßig
	mit dem Programm \texttt{MiKTeX Console} ein Update der installierten
	Pakete durchzuführen.} (\latex-Basisumgebung),
\item \textbf{TeXstudio}%
	\footnote{\url{https://www.texstudio.org/}}
	(Editor, unterstützt UTF-8 und beinhaltet einen integrierten PDF-Viewer).
\end{enumerate}
%
Alternative Editoren und PDF-Viewer:
%
\begin{enumerate}
	\item Visual Studio Code%
	\footnote{\url{https://code.visualstudio.com/}}
	mit LaTeX Workshop Extension,%
	\footnote{\url{https://marketplace.visualstudio.com/items?itemName=James-Yu.latex-workshop}}
	\item IntelliJ IDEA,%
	\footnote{\url{https://www.jetbrains.com/idea/}}
	mit TeXiFy-IDEA Plugin,%
	\footnote{\url{https://plugins.jetbrains.com/plugin/9473-texify-idea}}
	\item LyX,%
	\footnote{\url{https://www.lyx.org/}}
	\item TeXworks,%
	\footnote{\url{https://www.tug.org/texworks/}}
	\item WinEdt,%
	\footnote{\url{https://www.winedt.com/}}
	\item Sumatra PDF ("\latex-freundlicher" PDF-Viewer).%
	\footnote{\url{https://www.sumatrapdfreader.org/}}
\end{enumerate}

\subsection{Verwendung unter macOS}

Für macOS empfiehlt sich die folgende Konfiguration:
%
\begin{enumerate}
\item 
	\textbf{MacTeX}%
	\footnote{\url{https://tug.org/mactex/} -- Achtung: Aktuelle
	MacTeX-Distributionen verlangen in der Regel eine weitgehend aktuelle Version
	von macOS. Auf älteren Versionen kann alternativ \emph{TeX Live} mit einem
	speziellen Installationsscript installiert werden. Um die Pakete der
	\latex-Distribution aktuell zu halten, sollte regelmäßig das \emph{TeX Live
	Utility} ausgeführt werden.}
	(\latex-Basisumgebung),
\item \textbf{TeXstudio} (Editor, unterstützt UTF-8 und beinhaltet einen
integrierten PDF-Viewer).
\end{enumerate}
%
Alternative Editoren und PDF-Viewer:
%
\begin{enumerate}
	\item Visual Studio Code mit LaTeX Workshop Extension,%
	\item LyX,%
	\item TeXworks,%
	\item Skim (\latex-freundlicher PDF-Viewer).%
	\footnote{\url{https://skim-app.sourceforge.io/}}
\end{enumerate}


\subsection{Verwendung unter Linux}

Unter Linux kann folgendes Setup zum Einsatz kommen:
%
\begin{enumerate}
	\item 
	\textbf{TeX Live}%
	\footnote{\url{https://tug.org/texlive/} -- Eine Installation unter Linux
	erfolgt -- abhängig von der verwendeten Distribution -- am einfachsten mit
	Hilfe des jeweiligen Paketverwaltungssystems (\zB \texttt{apt-get}).}
	(\latex-Basisumgebung),
	\item \textbf{TeXstudio} (Editor, unterstützt UTF-8 und beinhaltet einen
	integrierten PDF-Viewer).
\end{enumerate}
%
Alternative Editoren und PDF-Viewer:
%
\begin{enumerate}
	\item Visual Studio Code mit LaTeX Workshop Extension,%
	\item LyX,%
	\item TeXworks,%
	\item qpdfview (\latex-freundlicher PDF-Viewer).%
	\footnote{\url{https://launchpad.net/qpdfview}}
\end{enumerate}

\subsection{Verwendung von Online-Umgebungen für \latex}

Neben einer lokalen \latex-Installation mit Editor gibt es mittlerweile auch
gute Online-Umgebungen, die das Erstellen von \latex-Dokumenten direkt im Browser
ermöglichen. Die \latex-Basisumgebung ist dabei auf den Servern des Dienstes
installiert. Dokumente können im Online-Editor erstellt oder auch bestehende
Vorlagen (wie etwa dieses Dokument) hochgeladen und weiter bearbeitet werden.
Die meisten Plattformen ermöglichen darüber hinaus ein kollaboratives
Arbeiten an einem Dokument.

Bei Verwendung solcher Umgebungen ist es sehr zu empfehlen, während
der Arbeit regelmäßig \emph{Backups} der Online Daten durchzuführen, um 
im schlimmsten Fall nicht wieder von vorne beginnen zu müssen.

\subsubsection{Overleaf}

Der bekannteste und mit dieser Vorlage getestete Editor ist \emph{Overleaf}%
\footnote{\url{https://www.overleaf.com/}}.
Um schnell Vorlagendokumente aus dem \texttt{hagenberg-thesis} Paket zu
importieren, können die Im\-port-Links im \emph{Readme}-Abschnitt zum
Github-Repository dieser Vorlage%
\footnote{\url{https://github.com/Digital-Media/HagenbergThesis}}
direkt verwendet werden.

Zu beachten gilt, dass kostenlose Overleaf Accounts seit Ende 2023
verkürzte Com\-pile-Timeouts von nur mehr 20 Sekunden aufweisen. Dies hat zur
Folge, dass größere Abschlussarbeiten sowie auch dieses Vorlagendokument damit
nicht mehr erstellt werden können, da die PDF-Generierung nach 20 Sekunden
abbricht.

Abhilfe schafft hier nur ein bezahlter Account (Lizenz) oder alternativ der
Verzicht auf Overleaf und der Umstieg auf ein lokales Setup.
In Bildungseinrichtungen, die über Lizenzen für Lehrende verfügen, können diese
Dokumente erstellen und zur Kollaboration an den*die Studierende*n freigeben, um
ein reibungsloses Arbeiten zu ermöglichen.

\subsubsection{Andere Online-Services}

Daneben gibt es noch weitere Online-Umgebungen für \latex\ und ihre Zahl wächst
beständig, \zB:
%
\begin{enumerate}
	\item Papeeria,%
	\footnote{\url{https://papeeria.com/}}
	\item CoCalc.%
	\footnote{\url{https://cocalc.ai/}}
\end{enumerate}
